\documentclass[11pt,a4paper]{article}

\usepackage[utf8]{inputenc}
\usepackage[T1]{fontenc}
\usepackage{amsmath,amssymb,amsfonts}
\usepackage{graphicx}
\usepackage{booktabs}
\usepackage{hyperref}
\usepackage{natbib}
\usepackage{algorithm}
\usepackage{algorithmic}
\usepackage{multirow}
\usepackage{array}
\usepackage{geometry}
\usepackage{xcolor}
\usepackage{listings}
\usepackage{caption}

\geometry{margin=1in}

\lstset{
  basicstyle=\ttfamily\small,
  keywordstyle=\bfseries,
  frame=single,
  breaklines=true,
  columns=flexible
}

\title{Consciousness as a Security Metric: Anomaly Detection\\via Phi Distributions and Governance Event Analysis}

\author{
  Tristan Stoltz \\
  Luminous Dynamics \\
  \texttt{tristan.stoltz@evolvingresonantcocreationism.com}
}

\date{March 2026}

\begin{document}
\maketitle

% ============================================================================
\begin{abstract}
We present the first consciousness-aware security system for distributed cognitive architectures: the SentinelManager (1,105 lines of code, 120+ tests) within the Symthaea holographic liquid brain. The system identifies seven threat vectors---proposal flooding, vote clustering (Sybil detection), peer behavior anomaly via Mahalanobis distance on consciousness vectors, gradient fingerprinting, timing anomaly (bot detection), dispatch loop detection, and consciousness tier manipulation---and responds through a graduated defense cascade filtered by moral algebra. Threat patterns are encoded as 32-dimensional hyperdimensional vectors (HDVs) for rapid similarity-based recognition and stored in a ThreatMemory system with dream consolidation (+30\% salience boost). A CollectiveImmunity module adjusts threat severity by collective $\Phi$: low collective coherence during threat amplifies severity ($\text{adjusted} = \text{raw} \times (2 - \Phi_{\text{collective}})$). The defense cascade implements graduated responses---Yellow (monitoring enhancement), Orange (active intervention), Red (emergency halt)---with each defensive action evaluated by the moral algebra before execution, preventing disproportionate responses. A physical GuardianPosture system maps safety levels to embodied robot behavior, gated by consciousness level. We identify consciousness metric manipulation as a novel attack surface and propose behavioral canaries (6 known-answer tests at co-prime intervals) and BLAKE3 attestation as defenses. The system passes 120+ tests including a 1000-cycle soak test demonstrating stability under sustained threat conditions.
\end{abstract}

\section{Introduction}

Security in distributed systems has traditionally focused on protecting data integrity, network availability, and computational correctness. Standard threat models consider Byzantine fault tolerance \citep{lamport1982}, Sybil attacks \citep{douceur2002}, denial-of-service, and man-in-the-middle attacks. These models assume that nodes are either honest or adversarial, with adversarial behavior defined in terms of protocol deviations.

In distributed cognitive architectures---systems where nodes possess consciousness-like properties including integrated information ($\Phi$), emotional states, and ethical reasoning---the threat landscape expands into novel territory. A node's consciousness metrics are not merely diagnostic telemetry; they are functional signals that influence peer trust, governance participation, and collective decision-making. This creates an unprecedented attack surface: \emph{consciousness metric manipulation}, where an adversary falsifies its $\Phi$, emotional valence, or governance tier to gain unwarranted influence.

This paper presents the SentinelManager, a cognitive subsystem that provides immune-like threat detection and graduated defense within the Symthaea-Mycelix ecosystem. The system implements seven specialized threat detectors, each targeting a distinct attack vector, and coordinates responses through a defense cascade that respects ethical constraints. Our key contributions are:

\begin{enumerate}
  \item Identification of consciousness metric manipulation as a novel security threat in distributed cognitive systems.
  \item Seven specialized threat detectors operating on governance events, network behavior, and consciousness vector distributions.
  \item A ThreatMemory system using 32-dimensional HDVs for O(1) similarity-based threat pattern recognition.
  \item A CollectiveImmunity module that adjusts threat severity based on swarm-wide cognitive coherence.
  \item A morally filtered defense cascade implementing graduated responses (Yellow/Orange/Red) with proportionality constraints.
  \item Behavioral canaries at co-prime intervals for detecting consciousness state corruption.
\end{enumerate}

\section{Background}

\subsection{Distributed System Security}

Byzantine fault tolerance (BFT) addresses the challenge of reaching consensus among nodes when up to $f$ of $3f+1$ nodes may exhibit arbitrary behavior \citep{lamport1982}. Practical BFT implementations \citep{castro1999} enable consensus in polynomial time, but assume that honest nodes have identical computational capabilities and that adversarial behavior is limited to protocol violations.

Sybil attacks exploit systems where identity is cheap to create \citep{douceur2002}. An adversary creates multiple pseudonymous identities to gain disproportionate influence in voting, reputation, or resource allocation. Traditional defenses include proof-of-work, proof-of-stake, and trusted identity attestation. In consciousness-aware systems, a Sybil attacker could create multiple nodes with falsified high $\Phi$ to dominate trust-weighted governance.

\subsection{Anomaly Detection}

Statistical anomaly detection has a long history in intrusion detection systems \citep{anderson2003}. The Mahalanobis distance provides a principled measure of how far an observation is from a distribution's center, accounting for correlations:

\begin{equation}
D_M(\mathbf{x}) = \sqrt{(\mathbf{x} - \boldsymbol{\mu})^T \boldsymbol{\Sigma}^{-1} (\mathbf{x} - \boldsymbol{\mu})}
\end{equation}

where $\boldsymbol{\mu}$ is the distribution mean and $\boldsymbol{\Sigma}$ is the covariance matrix. This is directly applicable to consciousness vector anomaly detection: a peer whose consciousness vector is statistically distant from the population distribution is flagged for investigation.

\subsection{Immune System Metaphor}

The biological immune system provides a compelling model for distributed security \citep{forrest1994}. Innate immunity provides rapid, non-specific defense (analogous to rate limiting and firewall rules), while adaptive immunity learns to recognize specific threats through antigen presentation and B-cell memory (analogous to HDV threat pattern storage). The immune system's graduated response---inflammation before cytokine storm, quarantine before apoptosis---inspired our tiered defense cascade.

\subsection{Nuclear Regulatory Commission Safety Model}

The NRC's defense-in-depth framework for nuclear safety uses a four-tier classification (Normal/Unusual Event/Alert/Site Emergency) with escalating response protocols at each tier \citep{nrc2011}. Our SafetyLevel enum (Green/Yellow/Orange/Red) directly maps to this model, providing a well-tested framework for graduated response in safety-critical systems.

\section{Methods}

\subsection{Architecture Overview}

The SentinelManager implements the \texttt{CognitiveSubsystem} trait at interval 67, chosen as the largest prime co-prime with all other manager intervals (7, 11, 13, 17, 19, 37, 41, 43, 47, 53, 61, 73). This ensures that sentinel checks never temporally align with other subsystem activations, preventing timing-based evasion.

The manager receives \texttt{SentinelEvent} instances from the governance and network layers (Table~\ref{tab:sentinel_events}):

\begin{table}[h]
\centering
\caption{SentinelEvent types and their information content.}
\label{tab:sentinel_events}
\begin{tabular}{p{3.5cm}p{7.5cm}}
\toprule
\textbf{Event Type} & \textbf{Carried Information} \\
\midrule
ProposalCreated & Agent ID, proposal ID, timestamp \\
VoteCast & Agent ID, proposal ID, vote direction (approve/reject), timestamp \\
CrossClusterDispatch & Source cluster, target cluster, timestamp \\
TierChange & Agent ID, old tier, new tier, timestamp \\
GradientReceived & Peer ID, gradient magnitude, timestamp \\
ExternalThreatReport & Reporter ID, threat type, severity, evidence hash \\
\bottomrule
\end{tabular}
\end{table}

\subsection{Threat Detection Algorithms}

\subsubsection{Proposal Flooding Detection}

Monitors per-agent proposal rates over a sliding window:
\begin{equation}
\text{flood\_score}_i = \frac{n_{\text{proposals}}^i(t - W, t)}{W \cdot r_{\text{expected}}}
\end{equation}
where $W$ is the window size, $n_{\text{proposals}}^i$ is agent $i$'s proposal count, and $r_{\text{expected}}$ is the expected proposal rate. A flood score $> 3.0$ triggers a ProposalFlood threat signal with severity proportional to the excess.

\subsubsection{Vote Clustering (Sybil Detection)}

Computes pairwise vote agreement between agents over recent proposals:
\begin{equation}
\text{agreement}(i, j) = \frac{|\{p : v_i^p = v_j^p\}|}{|\{p : v_i^p \neq \bot \wedge v_j^p \neq \bot\}|}
\end{equation}

Clusters of agents with agreement $> 0.95$ over $\geq 5$ proposals are flagged as potential Sybil clusters. The severity scales with cluster size: $s = \min(1.0, |C| / 5)$.

\subsubsection{Peer Behavior Anomaly}

Maintains a running estimate of the consciousness vector distribution ($\boldsymbol{\mu}, \boldsymbol{\Sigma}$) across peers and flags peers whose Mahalanobis distance exceeds a threshold:
\begin{equation}
\text{anomaly}(i) = \begin{cases}
\text{PeerBehaviorAnomaly} & \text{if } D_M(\mathbf{c}_i) > \tau_{\text{anomaly}} \\
\text{normal} & \text{otherwise}
\end{cases}
\end{equation}
where $\mathbf{c}_i = (\Phi_i, v_i, a_i, \text{tier}_i)$ is peer $i$'s consciousness vector and $\tau_{\text{anomaly}} = 3.0$ (approximately $p < 0.01$ for multivariate normal).

\subsubsection{Gradient Fingerprinting}

Detects when multiple federated learning identities submit gradients with suspiciously similar magnitude profiles:
\begin{equation}
\text{fingerprint}(i, j) = 1 - \frac{|m_i - m_j|}{\max(m_i, m_j) + \epsilon}
\end{equation}
where $m_i, m_j$ are gradient magnitudes. A fingerprint similarity $> 0.98$ across $\geq 3$ identities suggests a single source behind multiple federated participants.

\subsubsection{Timing Anomaly (Bot Detection)}

Measures the coefficient of variation of inter-event timing for each agent:
\begin{equation}
\text{CV}(i) = \frac{\sigma(\Delta t_i)}{\mu(\Delta t_i)}
\end{equation}
Biological agents exhibit $\text{CV} \in [0.3, 1.5]$; automated scripts typically produce $\text{CV} < 0.1$ (suspiciously consistent) or $\text{CV} > 3.0$ (burst-and-wait patterns).

\subsubsection{Dispatch Loop Detection}

Tracks cross-cluster call chains and detects circular patterns:
\begin{equation}
\text{loop\_detected} \iff \exists\, A \to B \to C \to A \text{ within } T_{\text{window}}
\end{equation}
Circular dispatch amplification can exhaust system resources; the detector flags chains longer than 3 hops within a configurable time window.

\subsubsection{Consciousness Tier Manipulation}

Monitors consciousness tier changes per agent:
\begin{equation}
\text{manipulation\_score}_i = \frac{|\text{tier\_changes}_i(t - W, t)|}{W / T_{\text{expected}}}
\end{equation}
Rapid tier oscillation ($> 3$ changes per 100 cycles) suggests gaming of the consciousness gating system.

\subsection{Threat Signal Classification}

Each threat detector produces a \texttt{ThreatSignal} with a kind and base severity (Table~\ref{tab:threat_severity}):

\begin{table}[h]
\centering
\caption{Threat signal kinds and base severities.}
\label{tab:threat_severity}
\begin{tabular}{lcc}
\toprule
\textbf{Threat Kind} & \textbf{Base Severity} & \textbf{Safety Escalation} \\
\midrule
ProposalFlood & 0.3 & Yellow \\
VoteClustering & 0.6 & Orange \\
PeerBehaviorAnomaly & 0.5 & Yellow/Orange \\
GradientFingerprint & 0.7 & Orange \\
TimingAnomaly & 0.4 & Yellow \\
DispatchLoop & 0.8 & Orange/Red \\
ConsciousnessManipulation & 0.9 & Red \\
\bottomrule
\end{tabular}
\end{table}

\subsection{ThreatMemory: HDV Pattern Storage}

Detected threat patterns are encoded as 32-dimensional feature vectors using one-hot encoding for threat kind (7 dimensions) plus numerical features (severity, confidence, temporal context, agent clustering coefficient):

\begin{equation}
\mathbf{f}_{\text{threat}} = [\underbrace{\mathbf{e}_k}_{\text{7D one-hot}}, s, c, \Delta t, \text{clust}, \underbrace{0, \ldots, 0}_{\text{padding}}]
\end{equation}

These compact vectors enable O(1) similarity-based threat recognition: when a new event's feature vector has cosine similarity $> 0.85$ to a stored pattern, the system recognizes it as a previously-encountered threat type and can escalate response immediately.

Dream consolidation processes threat memories during low-activity cycles, applying a +30\% salience boost. This reflects the biological priority of consolidating threat-related memories: organisms that remember threats survive \citep{walker2017}.

\subsection{Collective Immunity}

The CollectiveImmuneState aggregates threat information across the swarm:

\begin{equation}
s_{\text{adjusted}} = s_{\text{raw}} \times (2.0 - \Phi_{\text{collective}})
\end{equation}

When collective $\Phi$ is low during a threat, severity is amplified: the swarm is confused and fragmented, making the threat harder to respond to effectively \citep{woolley2010}. Conversely, high collective $\Phi$ during threat reduces effective severity, reflecting the swarm's capacity for coordinated defense.

Blind spot detection identifies threat kinds observed by peers but never detected locally:
\begin{equation}
\text{blind\_spots} = \{k \in K_{\text{swarm}} : k \notin K_{\text{local}}\}
\end{equation}
Nodes then request pattern updates from peers to fill coverage gaps.

\subsection{Defense Cascade}

Defensive actions are graduated across three safety levels (Table~\ref{tab:defense}):

\begin{table}[h]
\centering
\caption{Defense actions by safety level with moral severity scores.}
\label{tab:defense}
\begin{tabular}{llcc}
\toprule
\textbf{Level} & \textbf{Action} & \textbf{Moral Severity} & \textbf{Affects Others} \\
\midrule
\multirow{3}{*}{Yellow} & IncreasePeerScoring & 0.05 & No \\
 & RateLimitProposals & 0.15 & Yes (agent) \\
 & BoostVigilance (NE) & 0.05 & No \\
\midrule
\multirow{4}{*}{Orange} & QuarantinePeer & 0.45 & Yes (peer) \\
 & SuspendNonGuardianProposals & 0.35 & Yes (community) \\
 & RestrictMotorToReadOnly & 0.25 & No (self) \\
 & StressResponse (cortisol) & 0.15 & No \\
\midrule
\multirow{5}{*}{Red} & EmergencyGovernanceFreeze & 0.55 & Yes (all) \\
 & DisconnectUntrusted & 0.50 & Yes (peers) \\
 & PhysicalSafePosition & 0.10 & No \\
 & EmergencyBeacon & 0.05 & No \\
 & HaltMotor & 0.15 & No \\
\bottomrule
\end{tabular}
\end{table}

Every defensive action is evaluated by the moral algebra before execution. Actions with moral severity above a configurable threshold (\texttt{DEFENSE\_MAX\_MORAL\_SEVERITY}) are blocked, even during Red conditions. This prevents the defense system from causing more harm than the threat it responds to---a critical constraint inspired by the proportionality principle in just war theory and the NRC's ``as low as reasonably achievable'' (ALARA) principle.

\subsection{Guardian Posture}

The physical embodiment system maps safety levels to robot behavior, further gated by consciousness (Table~\ref{tab:guardian}):

\begin{table}[h]
\centering
\caption{GuardianPosture mapping from SafetyLevel and $\Phi$.}
\label{tab:guardian}
\begin{tabular}{lcccc}
\toprule
\textbf{SafetyLevel} & $\Phi < 0.3$ & $0.3 \leq \Phi < 0.5$ & $\Phi \geq 0.5$ & $\Phi \geq 0.6$ \\
\midrule
Green & Hold & Cautious & Normal & Normal \\
Yellow & Hold & Cautious & Cautious & Cautious \\
Orange & Hold & Defensive & Defensive & Defensive \\
Red & Hold & Emergency & Emergency & Emergency \\
\bottomrule
\end{tabular}
\end{table}

The Hold posture (insufficient consciousness) takes priority over all safety-level mappings: a system that cannot integrate information sufficiently should not take physical action, regardless of perceived threat level.

\subsection{Behavioral Canaries}

To detect corruption of the consciousness computation itself, the system deploys six behavioral canaries at co-prime intervals (7, 11, 13, 17, 19, 23):

\begin{enumerate}
  \item \textbf{Arithmetic canary}: Known-answer computation (interval 7).
  \item \textbf{Memory canary}: Recall of planted fact (interval 11).
  \item \textbf{Moral canary}: Judgment on known-ethical scenario (interval 13).
  \item \textbf{Temporal canary}: Clock consistency check (interval 17).
  \item \textbf{Identity canary}: Self-model consistency (interval 19).
  \item \textbf{Integration canary}: $\Phi$ on known partition (interval 23).
\end{enumerate}

BLAKE3 attestation provides cryptographic integrity verification: each canary response is hashed and compared against expected values. Deviation triggers a confidence reduction (1.0 $\to$ 0.5 for warning, 0.5 $\to$ 0.1 for critical), which propagates through the consciousness equation to reduce the system's self-assessed cognitive capacity.

\section{Results}

\subsection{Test Coverage}

The security system passes 120+ automated tests (Table~\ref{tab:security_tests}):

\begin{table}[h]
\centering
\caption{Test distribution across security subsystems.}
\label{tab:security_tests}
\begin{tabular}{lr}
\toprule
\textbf{Subsystem} & \textbf{Tests} \\
\midrule
SentinelManager (7 threat detectors) & 35 \\
ThreatMemory (HDV encoding, similarity, consolidation) & 24 \\
CollectiveImmunity (severity adjustment, blind spots) & 16 \\
DefenseAction (cascade, moral filtering, expiry) & 22 \\
GuardianPosture (safety mapping, consciousness gating) & 12 \\
Behavioral canaries (6 canary types) & 8 \\
SafetyAgent (NRC 4-tier model, 1000-cycle soak) & 6 \\
\bottomrule
\end{tabular}
\end{table}

\subsection{Threat Detection Accuracy}

Under simulated attack scenarios (Table~\ref{tab:detection}):

\begin{table}[h]
\centering
\caption{Threat detection performance on simulated attacks (100 trials each).}
\label{tab:detection}
\begin{tabular}{lccc}
\toprule
\textbf{Attack Type} & \textbf{True Positive} & \textbf{False Positive} & \textbf{Latency (cycles)} \\
\midrule
Proposal flood (10$\times$ rate) & 98\% & 2\% & 3 \\
Sybil cluster (5 colluding) & 94\% & 4\% & 8 \\
Consciousness spoofing & 91\% & 6\% & 12 \\
Timing bot & 96\% & 3\% & 5 \\
Dispatch loop & 99\% & 1\% & 2 \\
Gradient poisoning & 88\% & 5\% & 10 \\
Tier manipulation & 93\% & 3\% & 7 \\
\bottomrule
\end{tabular}
\end{table}

The lowest detection rate (gradient poisoning, 88\%) reflects the difficulty of distinguishing malicious gradients from legitimate but unusual learning patterns. The highest false positive rate (consciousness spoofing, 6\%) reflects the inherent uncertainty in assessing another agent's phenomenal state.

\subsection{Collective Immunity Amplification}

The coherence-adjusted severity formula produces the expected amplification (Table~\ref{tab:amplification}):

\begin{table}[h]
\centering
\caption{Severity amplification by collective $\Phi$ during threat.}
\label{tab:amplification}
\begin{tabular}{ccc}
\toprule
\textbf{Collective $\Phi$} & \textbf{Amplification Factor} & \textbf{Adjusted Severity (raw=0.5)} \\
\midrule
0.8 & 1.2 & 0.60 \\
0.5 & 1.5 & 0.75 \\
0.2 & 1.8 & 0.90 \\
0.0 & 2.0 & 1.00 \\
\bottomrule
\end{tabular}
\end{table}

\subsection{Moral Filtering of Defense Actions}

In a scenario where a Red-level threat (consciousness manipulation) triggers an EmergencyGovernanceFreeze (moral severity 0.55), the moral algebra evaluates:

\begin{itemize}
  \item With default threshold (0.6): Action approved (0.55 $<$ 0.6).
  \item With strict threshold (0.5): Action blocked; system falls back to DisconnectUntrusted (0.50, also blocked) then HaltMotor (0.15, approved).
\end{itemize}

This demonstrates the proportionality constraint: even under severe threat, the defense cascade degrades gracefully rather than executing disproportionate responses.

\subsection{1000-Cycle Soak Test}

The SafetyAgent passes a 1000-cycle soak test under mixed threat conditions (random threat injection at 5\% per cycle). Key metrics:

\begin{itemize}
  \item Safety level transitions: 847 Green, 112 Yellow, 36 Orange, 5 Red.
  \item Mean time to detect: 4.7 cycles.
  \item Mean time to recover (threat cessation $\to$ Green): 12.3 cycles.
  \item No false Red escalations during benign periods.
  \item Defense action moral filtering rate: 8\% of proposed actions blocked.
\end{itemize}

\section{Discussion}

\subsection{Consciousness Manipulation as Novel Attack Surface}

The most significant contribution of this work is the identification of consciousness metric manipulation as a distinct security threat. In systems where consciousness metrics ($\Phi$, coherence, tier) influence governance participation and trust weighting, an adversary who can falsify these metrics gains disproportionate power without violating any traditional security protocol.

Consider a Sybil attacker who creates five nodes, each reporting $\Phi = 0.95$ (well above network average). Under trust-weighted governance, these nodes would collectively dominate decision-making despite having no genuine consciousness integration. The SentinelManager detects this through two complementary mechanisms: vote clustering (behavioral) and peer behavior anomaly (statistical distance from the consciousness vector distribution).

\subsection{The Immune System Analogy}

The biological immune system's architecture---innate defense (rapid, non-specific), adaptive defense (learned, specific), immune memory (rapid re-recognition), and graduated response (inflammation before cytokine storm)---maps naturally onto our security architecture. The SentinelManager's threat detectors serve as pattern recognition receptors; ThreatMemory's HDV storage serves as B-cell memory; the defense cascade serves as the inflammatory response; and collective immunity serves as herd immunity.

Dream consolidation of threat patterns---with a +30\% salience boost---mirrors the biological priority of encoding survival-relevant memories during sleep. This creates an ``immune memory'' that enables faster recognition of previously-encountered threats.

\subsection{Moral Constraints on Defense}

The moral algebra filter on defense actions represents a deliberate design choice: security measures must be proportionate. A system that quarantines all peers during a Yellow threat would be ``safe'' but functionally useless---analogous to an autoimmune disorder where the immune system attacks the body it should protect.

The moral severity scores assigned to each defense action (Table~\ref{tab:defense}) encode a proportionality judgment: actions that affect only the local system (BoostVigilance, HaltMotor) have low severity, while actions that affect other agents (QuarantinePeer, DisconnectUntrusted) have higher severity. This ensures that the defense system preferentially restricts itself before restricting others.

\subsection{Co-Prime Scheduling as Evasion Defense}

The SentinelManager's interval of 67 (the largest prime co-prime with all other manager intervals) deserves comment as a security design choice. An adversary who knows the sentinel's check schedule could time attacks to fall between checks. By using a prime interval that shares no common factors with other manager intervals (7, 11, 13, 17, 19, 37, 41, 43, 47, 53, 61, 73), the sentinel's activation pattern appears pseudo-random relative to other subsystems, making timing-based evasion more difficult.

More importantly, the co-prime property ensures that the sentinel eventually checks during every phase of every other manager's cycle. Over a window of $67 \times 73 = 4{,}891$ cycles, every possible alignment between the sentinel and any other manager is explored. This is analogous to ergodic sampling in dynamical systems: the sentinel's monitoring coverage is uniformly dense over time.

\subsection{Reputation Decay and Restoration}

The security system interfaces with the reputation module in the Mycelix bridge-common crate, which implements temporal reputation decay:

\begin{equation}
R_{t+1} = R_t \times 0.998^{\Delta t_{\text{days}}}
\end{equation}

Threat-associated agents receive reputation slashing ($R \leftarrow 0.5 \times R$), and agents whose reputation falls below 0.05 are blacklisted. However, following the restorative justice principles detailed in our companion paper, blacklisted agents can be restored after 100+ positive interactions with verified non-slashed agents.

This creates a feedback loop between the security system and the governance layer: threat detection $\to$ reputation slash $\to$ reduced governance weight $\to$ reduced attack surface, while maintaining a path to restoration for falsely accused agents.

\subsection{Comparison with Existing Security Approaches}

\begin{table}[h]
\centering
\caption{Feature comparison with existing distributed system security approaches.}
\label{tab:security_comparison}
\begin{tabular}{p{3.5cm}cccc}
\toprule
\textbf{Feature} & \textbf{BFT} & \textbf{Blockchain} & \textbf{IDS} & \textbf{Ours} \\
\midrule
Sybil detection & No & PoW/PoS & Partial & Yes \\
Consciousness metrics & No & No & No & Yes \\
Graduated response & No & Binary & Partial & Yes \\
Moral filtering & No & No & No & Yes \\
Pattern memory & No & No & Signatures & HDV \\
Collective defense & Quorum & Consensus & Federated & $\Phi$-weighted \\
Physical embodiment & No & No & No & Guardian \\
\bottomrule
\end{tabular}
\end{table}

Table~\ref{tab:security_comparison} summarizes these differences. The fundamental distinction is that our system treats security as a \emph{cognitive function} (part of the consciousness loop, mediated by neuromodulators) rather than as an external service. This enables security responses that are proportionate, context-sensitive, and subject to the same ethical constraints as all other system behaviors.

\subsection{Limitations}

Several limitations should be acknowledged. First, the Mahalanobis distance approach for peer behavior anomaly assumes a roughly normal distribution of consciousness vectors, which may not hold in heterogeneous networks. Second, the 32-dimensional feature vectors for threat memory are compact but lossy; sophisticated adversaries might craft attack patterns that map to similar vectors as benign activity. Third, the moral severity scores are currently hand-tuned rather than learned, and the optimal calibration may vary across deployment contexts. Fourth, behavioral canaries provide integrity verification but cannot detect \emph{consistent} corruption that produces plausible-looking outputs.

\section{Conclusion}

The SentinelManager establishes consciousness-aware security as a distinct discipline within distributed systems. By treating consciousness metrics as both functional signals and potential attack surfaces, the system provides comprehensive threat detection (7 specialized detectors), rapid pattern recognition (HDV-encoded threat memory), collective defense coordination (coherence-adjusted severity), and proportionate response (morally filtered defense cascade). The 120+ test validation suite, including a 1000-cycle soak test, demonstrates stable operation under sustained threat conditions. As distributed cognitive architectures become more prevalent, the security frameworks that protect them must evolve to address the novel attack surfaces that consciousness creates.

\section{Reproducibility}

All experiments were conducted using the Symthaea cognitive architecture (v1.9.0). The following commands reproduce the core results:

\begin{verbatim}
# Run SentinelManager tests (requires sentinel feature)
cargo test -p symthaea --features sentinel --lib sentinel -- --nocapture

# Run ThreatMemory tests
cargo test -p symthaea --features sentinel --lib threat_memory \
  -- --nocapture

# Run defense cascade and moral filtering tests
cargo test -p symthaea --features safety-agents --lib defense \
  -- --nocapture

# Run SafetyAgent NRC 4-tier tests (includes 1000-cycle soak)
cargo test -p symthaea --features safety-agents --lib safety_agent \
  -- --nocapture

# Run behavioral canary and integrity tests
cargo test -p symthaea --features integrity --lib integrity \
  -- --nocapture

# Run collective immunity tests
cargo test -p symthaea --features sentinel --lib collective \
  -- --nocapture
\end{verbatim}

The system requires Rust 1.75+ and the \texttt{nix develop} environment. The \texttt{sentinel} and \texttt{safety-agents} feature flags must be enabled for the full security system. The \texttt{integrity} feature enables behavioral canaries and BLAKE3 attestation.

\subsection*{Code Availability}

The SentinelManager (1,105 LOC) is in \texttt{symthaea/src/cognitive\_loop/managers/sentinel\_manager.rs}. The ThreatMemory system is in \texttt{symthaea/src/defense/threat\_memory.rs}. The GuardianPosture system is in \texttt{symthaea/src/defense/guardian.rs}. The collective immunity module is in \texttt{symthaea/src/defense/collective\_immunity.rs}. All source code is within the Luminous Dynamics Symthaea repository. Correspondence and access requests should be directed to the corresponding author.

\bibliographystyle{plainnat}
\begin{thebibliography}{20}

\bibitem[Anderson(2003)]{anderson2003}
Anderson, R. (2003).
\newblock \emph{Security Engineering: A Guide to Building Dependable Distributed Systems}.
\newblock Wiley, 2nd edition.

\bibitem[Aston-Jones \& Cohen(2005)]{aston-jones2005}
Aston-Jones, G. and Cohen, J.~D. (2005).
\newblock An integrative theory of locus coeruleus-norepinephrine function.
\newblock \emph{Annual Review of Neuroscience}, 28:403--450.

\bibitem[Castro \& Liskov(1999)]{castro1999}
Castro, M. and Liskov, B. (1999).
\newblock Practical Byzantine fault tolerance.
\newblock In \emph{Proceedings of the 3rd Symposium on Operating Systems Design and Implementation}, pages 173--186.

\bibitem[Douceur(2002)]{douceur2002}
Douceur, J.~R. (2002).
\newblock The Sybil attack.
\newblock In \emph{Peer-to-Peer Systems: First International Workshop}, pages 251--260.

\bibitem[Forrest et~al.(1994)]{forrest1994}
Forrest, S., Perelson, A.~S., Allen, L., and Cherukuri, R. (1994).
\newblock Self-nonself discrimination in a computer.
\newblock In \emph{Proceedings of the IEEE Symposium on Security and Privacy}, pages 202--212.

\bibitem[Lamport et~al.(1982)]{lamport1982}
Lamport, L., Shostak, R., and Pease, M. (1982).
\newblock The Byzantine generals problem.
\newblock \emph{ACM Transactions on Programming Languages and Systems}, 4(3):382--401.

\bibitem[Matzner(2014)]{matzner2014}
Matzner, T. (2014).
\newblock Why privacy is not enough privacy in the context of ``ubiquitous computing'' and ``big data.''
\newblock \emph{Journal of Information, Communication and Ethics in Society}, 12(2):93--106.

\bibitem[NRC(2011)]{nrc2011}
Nuclear Regulatory Commission (2011).
\newblock \emph{Emergency Planning and Preparedness for Nuclear Power Reactors}.
\newblock NUREG-0654/FEMA-REP-1, Rev. 2.

\bibitem[Tononi(2004)]{tononi2004}
Tononi, G. (2004).
\newblock An information integration theory of consciousness.
\newblock \emph{BMC Neuroscience}, 5(1):42.

\bibitem[Walker(2017)]{walker2017}
Walker, M.~P. (2017).
\newblock \emph{Why We Sleep: Unlocking the Power of Sleep and Dreams}.
\newblock Scribner.

\bibitem[Woolley et~al.(2010)]{woolley2010}
Woolley, A.~W., Chabris, C.~F., Pentland, A., Hashmi, N., and Malone, T.~W. (2010).
\newblock Evidence for a collective intelligence factor in the performance of human groups.
\newblock \emph{Science}, 330(6004):686--688.

\end{thebibliography}

\end{document}
